\section{The quasi-Albanese and the ordinary fundamental group}\label{sec:albanese}

The fibre description in Section~\ref{sec:definitions} shows that
$a:X\to E$ has connected fibres. We now identify $a$ by its
quasi-Albanese universal property and compute the ordinary fundamental
group of $X$. Neither argument uses specialness.

Choose $q_*\notin\{q_i\}$. The unchanged complete curve
$E\times\{q_*\}$ gives a section
\begin{equation}\label{eq:section}
 s:E\longrightarrow X,
 \qquad a\circ s=\operatorname{id}_E.
\end{equation}

\subsection{Units and the universal property}
\begin{lemma}\label{lem:units}
Every invertible regular function on $X$ is constant.
\end{lemma}
\begin{proof}
For $v\in\Gamma(X,\OO_X)^*$, its rational extension to $Y$ has divisor
supported on the boundary, say $\divisor_Y(v)=\sum_i b_iD_i$. Intersect
with $A_j$ and use \eqref{eq:intersection}:
\[
 0=\divisor_Y(v)\cdot A_j=b_j.
\]
Thus the rational function has no poles anywhere on the normal
projective variety $Y$. It is regular and hence constant.
\end{proof}

Let $G$ be any semiabelian variety and let $f:X\to G$ be an algebraic
morphism. Write
\[
 1\longrightarrow T\longrightarrow G\xrightarrow{\rho}B
 \longrightarrow1,
 \qquad T\simeq\Gm^r,
\]
with $B$ abelian. For $e$ outside the finite set $\{e_i\}$, the fibre
$a^{-1}(e)$ is the punctured line $\PP^1\setminus\{q_i\}$. The
restriction of $\rho\circ f$ to this fibre extends to $\PP^1$ because
$B$ is proper. Every morphism $\PP^1\to B$ is constant: invariant
one-forms span the cotangent space of $B$ everywhere, and their pullbacks
to $\PP^1$ vanish. Its differential is zero, which implies constancy in
characteristic zero.

Comparing with the point supplied by the section $s$ gives
\[
 \rho\circ f=(\rho\circ f\circ s)\circ a
\]
on a dense open subset of $X$ and hence everywhere. The morphism
\[
 h(x)=f(x)-f(s(a(x)))
\]
therefore has image in $T$. Each coordinate of $h:X\to\Gm^r$ is a unit,
so Lemma~\ref{lem:units} makes $h$ constant. It vanishes on $s(E)$, and
hence vanishes identically. We have proved the unique factorization
\begin{equation}\label{eq:universal-factorization}
 f=(f\circ s)\circ a.
\end{equation}
Uniqueness follows by composing with $s$.

For the based universal property, choose the base point $s(0)$ and
suppose $f(s(0))=0$. Then $g=f\circ s:E\to G$ is a group homomorphism.
Indeed, its projection to $B$ is a homomorphism by the rigidity theorem
for abelian varieties. Its group-law defect
\[
 E\times E\longrightarrow G,
 \qquad(e,e')\longmapsto g(e+e')-g(e)-g(e')
\]
therefore takes values in the affine group $T$. A morphism from a
connected projective variety to an affine torus is constant, and the
defect is zero at $(0,0)$. This proves the claim. Together with
\eqref{eq:universal-factorization}, it establishes
\[
 \Alb(X)=E,\qquad a_X=a.
\]
The general fibre in \eqref{eq:general-a-fibre} has logarithmic
canonical degree $n-2>0$. Its identity fibration is of general type, so
it is not special. Together with connectedness of all fibres, this
proves part~(ii) of Theorem~\ref{thm:example}.

\begin{remark}[Logarithmic one-forms]\label{rem:log-forms}
The same intersection matrix gives a cohomological check. The
logarithmic residue sequence is
\[
 0\longrightarrow\Omega_Y^1\longrightarrow\Omega_Y^1(\log D)
 \longrightarrow\bigoplus_i\OO_{D_i}\longrightarrow0.
\]
The connecting map sends a constant residue vector $(c_i)$ to the
corresponding linear combination of divisor classes. The classes $[D_i]$
are independent, since pairing with $[A_j]$ gives $c_j$. Consequently
\[
 H^0(Y,\Omega_Y^1(\log D))=H^0(Y,\Omega_Y^1)
   =\bar a^*H^0(E,\Omega_E^1).
\]
The boundary creates no additional logarithmic one-form. The universal
property above also rules out an isogeny ambiguity without needing to
deduce the whole quasi-Albanese from this dimension count.
\end{remark}

\subsection{The meridians die inside the given open surface}
Remove the retained exceptional curves temporarily and write
\[
 U=X\setminus\bigcup_iA_i^\circ.
\]
The blowup is an isomorphism away from its centres, and their containing
horizontal fibres have already been deleted. Thus
\begin{equation}\label{eq:product-open}
 U\simeq E\times(\PP^1\setminus Q),
 \qquad Q=\{q_1,\ldots,q_n\}.
\end{equation}
Let $m_i$ denote positively oriented puncture loops in
$\PP^1\setminus Q$, with paths to a common base point. They generate its
fundamental group with the single relation $m_1\cdots m_n=1$.

Each $A_i^\circ$ is a closed smooth complex divisor in $X$. The inclusion
$U\hookrightarrow X$ induces a surjection on fundamental groups, whose
kernel is normally generated by meridians around these divisors. For
completeness, loops can be perturbed off a real codimension-two
submanifold. A null-homotopy can be made transverse to it, relative to
the boundary of its parameter disc. It then meets the divisors in
finitely many points. Removing small discs about those points expresses
the boundary loop as a product of conjugates of meridians. This argument
uses compactness of the parameter disc, not compactness of the divisors.

The meridians can be identified explicitly. Near $(e_i,q_i)$ choose local
coordinates $(x,t)$ on $S$ so the deleted horizontal curve is $t=0$.
The blowup chart
\begin{equation}\label{eq:blowup-chart}
 b(u,t)=(ut,t)
\end{equation}
contains the whole retained affine line $A_i^\circ=\{t=0\}$, with
coordinate $u$. The strict transform of the deleted horizontal curve
does not meet this chart. At $u=0$ the loop
\[
 t=\varepsilon e^{2\pi i\theta},\qquad0\leq\theta\leq1,
\]
is a meridian of $A_i^\circ$. Under \eqref{eq:product-open}, it has fixed
elliptic coordinate $e_i$ and winds once about $q_i$. It represents a
conjugate of $m_i$. Moreover, it bounds the actual disc
$\{u=0,\ |t|\leq\varepsilon\}$ inside $X$, whose centre lies on
$A_i^\circ$.

It follows that
\begin{align*}
 \pi_1(X)
 &\simeq
 \frac{\pi_1(E)\times\pi_1(\PP^1\setminus Q)}
      {\langle\!\langle m_1,\ldots,m_n\rangle\!\rangle}\\
 &\simeq\pi_1(E)\simeq\Z^2.
\end{align*}
Under the product identification \eqref{eq:product-open}, $a$ is the
projection to $E$, so $a_*$ is precisely the displayed isomorphism. The
section \eqref{eq:section} also shows directly that the elliptic
generators survive. For a general fibre, all generators of its free
fundamental group are puncture loops, so its image in $\pi_1(X)$ is trivial. This
proves part~(iii) and completes Theorem~\ref{thm:example}. Every relation
used here is realized by a homotopy in $X$; no additional orbifold
relation has been imposed on its ordinary group.
